% ---------------------------------------------------------------------
%  TFPT --- Red Team / Stress Test layer (Targets A--E)
%  A deliberately adversarial pass over the five load-bearing reductions.
%  Backed by verification/redteam/rt_A..rt_F + run_redteam.py.
% ---------------------------------------------------------------------
\documentclass[11pt,a4paper]{article}
\usepackage[T1]{fontenc}
\usepackage[utf8]{inputenc}
\usepackage[english]{babel}
\usepackage{lmodern}
\usepackage[a4paper,margin=2.3cm]{geometry}
\usepackage{amsmath,amssymb,amsthm}
\usepackage{mathtools}
\usepackage{booktabs,array,tabularx}
\usepackage{enumitem}
\usepackage{xcolor}
\usepackage[most]{tcolorbox}
\usepackage{microtype}
\usepackage[hidelinks,colorlinks=true,linkcolor=blue!55!black,urlcolor=blue!55!black]{hyperref}
\emergencystretch=3em

% ---------------------------------------------------------------------
%  tfpt_style.tex -- THE shared visual style for every TFPT document.
%  Single source for colours, status markers, marker key, the \veri
%  script reference, the tcolorbox families, and the running
%  header/footer.  \input AFTER xcolor + tcolorbox[most] + hyperref are
%  loaded (every paper loads those).  Each document sets its short
%  running title before \begin{document}, e.g.
%        \def\TFPTdoctitle{Architecture and the E8 Compiler}
%  ---------------------------------------------------------------------
\usepackage{fancyhdr}
\usepackage{etoolbox}
\usepackage{tikz}
\usetikzlibrary{positioning,arrows.meta,fit,backgrounds,calc,shapes.geometric,shapes.misc}

% ---- single-source author / version / automatic date stamp ----
% ---------------------------------------------------------------------
%  tex-artefacts/version.tex -- single source for author list, version
%  and the automatic title-page date stamp of every TFPT document.
%  \input by tex-artefacts/tfpt_style.tex, so every paper that loads the
%  shared style automatically gets:
%     \TFPTauthors   -- the author list (use as \author{\TFPTauthors})
%     \TFPTversion   -- curated semantic version (bump per release; mirror
%                       the bump in changelog.tex)
%     \TFPTrev       -- auto build revision = git commit count, injected by
%                       build.sh into tex-artefacts/version-auto.tex
%                       (gitignored); falls back to "dev" for standalone runs
%     \TFPTdatestamp -- the title-page date line: \today + version + revision
%  build.sh also writes tex-artefacts/release-version.txt (gitignored):
%     "{TFPTversion} (rev {git commit count})" for Zenodo + release tooling.
%  Bump \TFPTversion here ONCE; it then updates in all papers at the next build.
% ---------------------------------------------------------------------
\def\TFPTauthors{Stefan Hamann \and Alessandro Rizzo}
\def\TFPTversion{5.3}
\def\TFPTrev{dev}
% build.sh regenerates the line below with the live git commit count:
\InputIfFileExists{tex-artefacts/version-auto.tex}{}{}
\def\TFPTdatestamp{\today\,\textperiodcentered\,v\TFPTversion\,(rev~\TFPTrev)}


% ---- colours (canonical TFPT palette) ----
\definecolor{tfblue}{HTML}{1F4E79}
\definecolor{tfgreen}{HTML}{1E7B45}
\definecolor{tfred}{HTML}{9B2226}
\definecolor{tfgold}{HTML}{B8860B}
\definecolor{tfgray}{HTML}{555555}

% ---- status markers (simplified 4-class display system) ----
% Reader-facing markers collapse to FOUR classes.  The fine-grained per-claim
% type (Axiom / Formal / Lattice / Numerical / Identity / Physical) stays in
% verification/status_ledger.csv -- the single source of truth -- so no fidelity
% is lost; the papers just show the coarse class so a reader is not overwhelmed.
%   [E] exact / machine-proven  (identity, Lie/lattice, formalised, numerical FP)
%   [C] conditional             (physical / bridge / readout, under named hypotheses)
%   [O] open / axiom            (declared input or genuine gap)
%   [X] kill test               (falsifiable threshold)
\newcommand{\mE}{{\small\textcolor{tfgreen!70!black}{\textbf{[E]}}}}
\newcommand{\mC}{{\small\textcolor{tfgold!78!black}{\textbf{[C]}}}}
\newcommand{\mO}{{\small\textcolor{tfred!80!black}{\textbf{[O]}}}}
\newcommand{\mX}{{\small\textcolor{tfred!58!black}{\textbf{[X]}}}}
% backward-compatible aliases: the old fine-grained names now render their class,
% so any not-yet-rewritten call site (or the standalone changelog) still compiles.
\newcommand{\tI}{\mE}\newcommand{\tL}{\mE}\newcommand{\tF}{\mE}\newcommand{\tN}{\mE}
\newcommand{\tIalg}{\mE}\newcommand{\tInum}{\mE}
\newcommand{\tP}{\mC}\newcommand{\tB}{\mC}\newcommand{\tR}{\mC}
\newcommand{\tA}{\mO}
\newcommand{\tX}{\mX}

% compact marker legend, placed under a marker-bearing table
\newcommand{\markerkey}{\par\vspace{1.5pt}\noindent{\scriptsize\itshape
  \textcolor{black!58}{Marker key:\ \mE~exact / machine-proven\,$\cdot$\,%
  \mC~conditional (holds under named hypotheses)\,$\cdot$\,%
  \mO~open / axiom\,$\cdot$\,\mX~falsifiable kill test.\ \
  Per-claim fine type in \texttt{verification/status\_ledger.csv}.}}\par\vspace{2pt}}

% horizontal badge bar (place once near the front of a document)
\newcommand{\TFPTbadges}{\par\noindent
  \fcolorbox{tfgreen!55!black}{tfgreen!8}{\scriptsize\mE\,exact / machine-proven}\hfill
  \fcolorbox{tfgold!70!black}{tfgold!10}{\scriptsize\mC\,conditional}\hfill
  \fcolorbox{tfred!75!black}{tfred!8}{\scriptsize\mO\,open / axiom}\hfill
  \fcolorbox{tfred!60!black}{tfred!6}{\scriptsize\mX\,kill test}\par\vspace{3pt}}

% inline reference to the script that machine-checks a claim (see verification/)
\newcommand{\veri}[1]{\,{\footnotesize\textcolor{tfblue!70!black}{[\,\texttt{verification/#1}\,]}}}
% lightweight inline colour-mark for a bare verification-script token in running
% prose / tables, e.g. \vref{v108}, \vref{v49}/\vref{v71}.  Same blue as \veri so a
% reader instantly sees "this is a machine-checked module reference".
\newcommand{\vref}[1]{\textcolor{tfblue!72!black}{\textsf{#1}}}

% ---- tcolorbox families (one visual language across all documents) ----
% keybox  : the central claim / reading of a section (blue)
% winbox  : an established result / closure (green)
% gapbox  : an open gate / honest caveat / kill criterion (red)
% auditbox: a recorded audit fingerprint, not load-bearing (gold)
% navbox  : a light, neutral navigation / reference card (document map, etc.)
% Visual language (deliberately light): a faint tint, a thin coloured frame
% and a light title band with dark coloured title text -- NOT a heavy
% white-on-saturated-colour bar -- so a page never reads as a wall of boxes.
\newtcolorbox{keybox}[1][]{enhanced,breakable,boxrule=0.5pt,arc=1.5pt,
  colback=tfblue!4!white,colframe=tfblue!70!black,coltitle=tfblue!88!black,colbacktitle=tfblue!12!white,
  fonttitle=\bfseries\small,fontupper=\small,title={#1},left=2mm,right=2mm,top=1mm,bottom=1mm}
% non-breaking variant (front-matter cards that should never split across a page)
\newtcolorbox{keyboxnb}[1][]{enhanced,breakable=false,boxrule=0.5pt,arc=1.5pt,
  colback=tfblue!4!white,colframe=tfblue!70!black,coltitle=tfblue!88!black,colbacktitle=tfblue!12!white,
  fonttitle=\bfseries\small,fontupper=\small,title={#1},left=2mm,right=2mm,top=1mm,bottom=1mm}
\newtcolorbox{winbox}[1][]{enhanced,breakable,boxrule=0.5pt,arc=1.5pt,
  colback=tfgreen!5!white,colframe=tfgreen!65!black,coltitle=tfgreen!50!black,colbacktitle=tfgreen!14!white,
  fonttitle=\bfseries\small,fontupper=\small,title={#1},left=2mm,right=2mm,top=1mm,bottom=1mm}
\newtcolorbox{gapbox}[1][]{enhanced,breakable,boxrule=0.5pt,arc=1.5pt,
  colback=tfred!4!white,colframe=tfred!70!black,coltitle=tfred!75!black,colbacktitle=tfred!12!white,
  fonttitle=\bfseries\small,fontupper=\small,title={#1},left=2mm,right=2mm,top=1mm,bottom=1mm}
% non-breaking variants (callouts that must stay whole on one page, embedded in text)
\newtcolorbox{winboxnb}[1][]{enhanced,breakable=false,boxrule=0.5pt,arc=1.5pt,
  colback=tfgreen!5!white,colframe=tfgreen!65!black,coltitle=tfgreen!50!black,colbacktitle=tfgreen!14!white,
  fonttitle=\bfseries\small,fontupper=\small,title={#1},left=2mm,right=2mm,top=1mm,bottom=1mm}
\newtcolorbox{gapboxnb}[1][]{enhanced,breakable=false,boxrule=0.5pt,arc=1.5pt,
  colback=tfred!4!white,colframe=tfred!70!black,coltitle=tfred!75!black,colbacktitle=tfred!12!white,
  fonttitle=\bfseries\small,fontupper=\small,title={#1},left=2mm,right=2mm,top=1mm,bottom=1mm}
\newtcolorbox{auditbox}[1][]{enhanced,breakable,boxrule=0.5pt,arc=1.5pt,
  colback=tfgold!7!white,colframe=tfgold!70!black,coltitle=tfgold!45!black,colbacktitle=tfgold!16!white,
  fonttitle=\bfseries\small,fontupper=\small,title={#1},left=2mm,right=2mm,top=1mm,bottom=1mm}
\newtcolorbox{navbox}[1][]{enhanced,breakable,boxrule=0.5pt,arc=1.5pt,
  colback=tfgray!3!white,colframe=tfgray!55,coltitle=tfgray!30!black,colbacktitle=tfgray!12!white,
  fonttitle=\bfseries\small,fontupper=\small,title={#1},left=2mm,right=2mm,top=1mm,bottom=1mm}
% back-compatibility aliases (older box names map onto the canonical types)
\newtcolorbox{audbox}[1][]{enhanced,breakable,boxrule=0.5pt,arc=1.5pt,
  colback=tfgold!7!white,colframe=tfgold!70!black,coltitle=tfgold!45!black,colbacktitle=tfgold!16!white,
  fonttitle=\bfseries\small,fontupper=\small,title={#1},left=2mm,right=2mm,top=1mm,bottom=1mm}
\newtcolorbox{resbox}[1][]{enhanced,breakable,boxrule=0.5pt,arc=1.5pt,
  colback=tfgreen!5!white,colframe=tfgreen!55!black,coltitle=tfgreen!50!black,colbacktitle=tfgreen!12!white,
  fonttitle=\bfseries\small,fontupper=\small,title={#1},left=2mm,right=2mm,top=1mm,bottom=1mm}
\newtcolorbox{roadbox}[1][]{enhanced,breakable,boxrule=0.5pt,arc=1.5pt,
  colback=tfgreen!5!white,colframe=tfgreen!55!black,coltitle=tfgreen!50!black,colbacktitle=tfgreen!12!white,
  fonttitle=\bfseries\small,fontupper=\small,title={#1},left=2mm,right=2mm,top=1mm,bottom=1mm}
\newtcolorbox{intbox}[1][]{enhanced,breakable,boxrule=0.5pt,arc=1.5pt,
  colback=tfgold!7!white,colframe=tfgold!70!black,coltitle=tfgold!45!black,colbacktitle=tfgold!16!white,
  fonttitle=\bfseries\small,fontupper=\small,title={#1},left=2mm,right=2mm,top=1mm,bottom=1mm}

% ---- colour-marked display formula (load-bearing key equations) ----
% A highlighted formula line: a coloured left accent bar + faint tint, no full
% frame -- so a key equation is visually set apart from the running prose
% without becoming yet another box.  Use:  \begin{keyeq}$\displaystyle ...$\end{keyeq}
\newtcolorbox{keyeq}{blanker,breakable=false,halign=center,
  colback=tfgold!8!white,borderline west={2.4pt}{0pt}{tfgold!75!black},
  left=8pt,right=6pt,top=3pt,bottom=3pt,before skip=5pt,after skip=5pt}

% ---- running header / footer (consistent across all documents) ----
\providecommand{\TFPTdoctitle}{}
\setlength{\headheight}{14pt}
\pagestyle{fancy}
\fancyhf{}
\fancyhead[L]{\footnotesize\textcolor{tfgray}{\textit{TFPT}\,\textbullet\,\TFPTdoctitle}}
\fancyhead[R]{\footnotesize\textcolor{tfgray}{\thepage}}
\fancyfoot[L]{\scriptsize\textcolor{tfgray}{v\TFPTversion\,\textperiodcentered\,rev~\TFPTrev}}
\renewcommand{\headrulewidth}{0.3pt}
\renewcommand{\footrulewidth}{0pt}
\fancypagestyle{plain}{\fancyhf{}\renewcommand{\headrulewidth}{0pt}%
  \fancyfoot[C]{\footnotesize\textcolor{tfgray}{\thepage}}%
  \fancyfoot[L]{\scriptsize\textcolor{tfgray}{v\TFPTversion\,\textperiodcentered\,rev~\TFPTrev}}}

% ---- shared figure macros (claim stack, anchor cockpit, E8 glue, 2/3 motif, K5) ----
%  (this fragment lives in tex-artefacts/ and is \input by documents compiled from
%   the repository root, so the path is given relative to that root.)
% ---------------------------------------------------------------------
%  tfpt_figures.tex -- shared, reusable TikZ figure macros for every
%  TFPT document.  \input by tfpt_style.tex (which loads tikz + the
%  needed libraries + etoolbox).  Each macro draws a self-contained
%  centred figure; nothing is drawn until the macro is called.
%
%    \TFPTclaimstack[zone]  -- the architecture claim stack (B3); the
%                              optional zone in {axioms,anchor,compiler,
%                              sm,bridges,frontier} is highlighted.
%    \TFPTanchorcockpit      -- the anchor a=(1,1,2) cockpit (B4)
%    \TFPTeightglue          -- the D5 (+) A3 +mu4 => E8 glue (B5)
%    \TFPTmotif              -- the 2/3 motif map (B6)
%    \TFPTactionkfive        -- the K5 action-lattice mnemonic (B7)
% ---------------------------------------------------------------------

% ===== B3: the architecture claim stack =========================
\newcommand{\tfhl}[1]{\ifdefstring{\tfcur}{#1}{cbhi}{cbnl}}
\newcommand{\TFPTclaimstack}[1][none]{%
\def\tfcur{#1}%
\begin{center}
\begin{tikzpicture}[font=\footnotesize,
  cb/.style={rounded corners=2pt,inner sep=4pt,text width=118mm,align=center,minimum height=8.5mm},
  cbnl/.style={line width=0.5pt}, cbhi/.style={line width=1.5pt},
  ar/.style={-{Stealth[length=2.2mm]},tfgray,semithick}]
  \node[cb,fill=tfred!8,draw=tfred!75!black,\tfhl{axioms}] (axioms)
    {\textbf{P1}\,$c_3=\tfrac1{8\pi}$ \quad \textbf{P2}\,$g_{\mathrm{car}}=5$ \ \textcolor{tfgray}{\itshape(two inputs, [O] axioms)}};
  \node[cb,fill=tfblue!8,draw=tfblue,\tfhl{anchor},below=3.5mm of axioms] (anchor)
    {Anchor $a=(1,1,2)$:\ $e_1,e_2,e_3=4,5,2$ \ \textcolor{tfgray}{\itshape([E] exact)}};
  \node[cb,fill=tfgreen!10,draw=tfgreen!70!black,\tfhl{compiler},below=3.5mm of anchor] (compiler)
    {Discrete compiler:\ $D_5\oplus A_3+\mu_4\Rightarrow E_8$,\ $h=30$,\ $\varphi_0$ \ \textcolor{tfgray}{\itshape([E] exact)}};
  \node[cb,fill=tfgreen!7,draw=tfgreen!60!black,\tfhl{sm},below=3.5mm of compiler] (sm)
    {SM skeleton:\ three families, hypercharges, flavor matrix $R$ \ \textcolor{tfgray}{\itshape([E])}};
  \node[cb,fill=tfgold!12,draw=tfgold!80!black,\tfhl{bridges},below=3.5mm of sm] (bridges)
    {Physical bridges:\ $v_{\mathrm{geo}}$,\ $G_{\mathrm{net}}$,\ $F_{\mathrm{transfer}}$ \ \textcolor{tfgray}{\itshape([C] bridges)}};
  \node[cb,fill=tfgray!12,draw=tfgray!70!black,\tfhl{frontier},below=3.5mm of bridges] (frontier)
    {Frontier readouts:\ $\Lambda$, inflation, $\eta_B$, dark matter, QG \ \textcolor{tfgray}{\itshape([C]/[O])}};
  \draw[ar] (axioms)--(anchor); \draw[ar] (anchor)--(compiler);
  \draw[ar] (compiler)--(sm);   \draw[ar] (sm)--(bridges); \draw[ar] (bridges)--(frontier);
\end{tikzpicture}
\end{center}}

% ===== B4: the anchor cockpit ====================================
\newcommand{\TFPTanchorcockpit}{%
\begin{center}
\begin{tikzpicture}[font=\footnotesize,
  pan/.style={rounded corners=2pt,draw=tfblue!60,fill=tfblue!4,inner sep=5pt,
    text width=46mm,align=left},
  hd/.style={font=\bfseries\small,tfblue}]
  \node[pan] (e) {\textbf{elementary} $e_k(a)$\\[2pt]
     $e_1=4\to|\mu_4|$\\ $e_2=5\to g_{\mathrm{car}}$\\ $e_3=2\to|\mathbb Z_2|$\\[2pt]
     $c_3=\tfrac1{2e_1\pi}=\tfrac1{8\pi}$};
  \node[pan,right=4mm of e] (p) {\textbf{power sums} $p_n=2+2^n$\\[2pt]
     $p_0,p_1,p_2,p_3=3,4,6,10$\\ $p_1p_2p_3=240=|R(E_8)|$\\ $p_4-p_3=8=\operatorname{rank}E_8$\\
     $\Rightarrow\dim E_8=248$};
  \node[pan,right=4mm of p] (m) {\textbf{derived motifs}\\[2pt]
     $4,5,8,16,25,$\\ $30,41,120,240$\\[2pt]
     \textcolor{tfgray}{\itshape one anchor,}\\ \textcolor{tfgray}{\itshape many projections}};
  \node[hd,above=1mm of p.north] {$a=(1,1,2)$ \ \textcolor{tfgray}{\footnotesize\itshape[E]}};
\end{tikzpicture}
\end{center}}

% ===== B5: the E8 glue ===========================================
\newcommand{\TFPTeightglue}{%
\begin{center}
\begin{tikzpicture}[font=\footnotesize,
  d/.style={rounded corners=2pt,draw=tfblue,fill=tfblue!6,inner sep=4pt,align=center,text width=36mm},
  mid/.style={rounded corners=2pt,draw=tfgold!80!black,fill=tfgold!12,inner sep=4pt,align=center,text width=52mm},
  e8/.style={rounded corners=2pt,draw=tfgreen!70!black,fill=tfgreen!10,inner sep=4pt,align=center,text width=52mm,font=\bfseries},
  rt/.style={rounded corners=2pt,draw=tfgreen!55!black,fill=tfgreen!6,inner sep=3pt,align=center,text width=52mm},
  side/.style={rounded corners=2pt,draw=tfgray!55,fill=tfgray!6,inner sep=4pt,align=left,text width=34mm,font=\scriptsize},
  ar/.style={-{Stealth[length=2mm]},tfgray,semithick}]
  \node[d] (d5) {$D_5$ discriminant\\ $\mathbb Z_4$, $q=5/4$};
  \node[d,right=10mm of d5] (a3) {$A_3$ discriminant\\ $\mathbb Z_4$, $q=3/4$};
  \node[mid,below=5mm of $(d5)!0.5!(a3)$] (glue) {$\mu_4$ anti-isometry\ ($q$-sum $=2$)};
  \node[e8,below=4mm of glue] (e8) {even unimodular $E_8$ lattice};
  \node[rt,below=4mm of e8] (rt) {$240$ roots,\ rank $8$,\ $\dim=248$};
  \node[side,right=8mm of glue] (why) {\textbf{why unique?}\\ rank $8$\\ even\\ unimodular\\ positive definite};
  \draw[ar] (d5)--(glue); \draw[ar] (a3)--(glue);
  \draw[ar] (glue)--(e8); \draw[ar] (e8)--(rt);
\end{tikzpicture}\\[2pt]
{\scriptsize\itshape $D_5\oplus A_3+\mu_4\Rightarrow E_8$ --- the unique even unimodular rank-$8$ glue \ [E] exact (lattice).}
\end{center}}

% ===== B5b: the double pushout (one glue, two categories) ========
\newcommand{\TFPTdoublepushout}{%
\begin{center}
\begin{tikzpicture}[font=\footnotesize,>={Stealth[length=2mm]},
  lat/.style={rounded corners=2pt,draw=tfblue,fill=tfblue!6,inner sep=5pt,align=center,text width=34mm},
  net/.style={rounded corners=2pt,draw=tfblue!70!black,fill=tfblue!4,inner sep=5pt,align=center,text width=34mm},
  e8l/.style={rounded corners=2pt,draw=tfgold!80!black,fill=tfgold!12,inner sep=5pt,align=center,text width=34mm,font=\bfseries},
  e8n/.style={rounded corners=2pt,draw=tfgreen!70!black,fill=tfgreen!10,inner sep=5pt,align=center,text width=34mm,font=\bfseries},
  ar/.style={->,tfgray,semithick},
  lb/.style={font=\scriptsize,text=black,align=center,fill=white,inner sep=1.5pt}]
  \node[lat] (LL) at (0,2.5)   {$D_5\oplus A_3$\\[1pt]\scriptsize root lattices,\ $\mathrm{disc}=\Z_4$};
  \node[e8l] (LR) at (6.6,2.5) {$E_8$\\[1pt]\scriptsize even unimodular,\ $\det 1$};
  \node[net] (BL) at (0,0)     {$(D_5)_1\otimes(A_3)_1$\\[1pt]\scriptsize chiral nets,\ $c=8$};
  \node[e8n] (BR) at (6.6,0)   {$(E_8)_1$\\[1pt]\scriptsize holomorphic,\ $\mu$-index $1$};
  \draw[ar] (LL)--node[lb]{$+\,\mu_4$ glue}(LR);
  \draw[ar] (BL)--node[lb]{$\rtimes\,\mu_4$ extension}(BR);
  \draw[ar] (LL)--node[lb,left=-1pt]{level-$1$}(BL);
  \draw[ar] (LR)--node[lb,right=-1pt]{level-$1$}(BR);
\end{tikzpicture}\\[2pt]
{\scriptsize\itshape One operation, two categories: the order-$|\mu_4|{=}4$ glue is the lattice pushout
$D_5\oplus A_3+\mu_4\Rightarrow E_8$ \emph{and} the simple-current / anyon-condensation extension
$(D_5)_1\otimes(A_3)_1\rtimes\mu_4\cong(E_8)_1$; the square commutes. \ [E] exact.}
\end{center}}

% ===== B6: the 2/3 motif map =====================================
\newcommand{\TFPTmotif}{%
\begin{center}
\begin{tikzpicture}[font=\footnotesize,
  ex/.style={rounded corners=2pt,draw=tfgreen!70!black,fill=tfgreen!8,inner sep=4pt,align=center,text width=92mm},
  br/.style={rounded corners=2pt,draw=tfgold!80!black,fill=tfgold!12,inner sep=4pt,align=center,text width=92mm},
  ar/.style={-{Stealth[length=2mm]},tfgray,semithick}]
  \node[ex] (a) {$|\mathbb Z_2|/N_{\mathrm{fam}}=2/3$ \ \textcolor{tfgray}{\itshape(anchor ratio)}};
  \node[ex,below=3mm of a] (b) {parabolic weights $\{0,\tfrac13,\tfrac23\}=\operatorname{spec}A_0^\star$};
  \node[ex,below=3mm of b] (c) {transfer spectrum $\{1,(2/3)^6,(1/3)^6\}$,\ gap $6\log\tfrac32$};
  \node[br,below=3mm of c] (d) {Koide target $2/3$ \ \textcolor{tfgray}{\itshape(source$\to$pole bridge [C])}};
  \node[br,below=3mm of d] (e) {Nariai / horizon entropy bound $2/3$ \ \textcolor{tfgray}{\itshape([C])}};
  \draw[ar] (a)--(b); \draw[ar] (b)--(c); \draw[ar] (c)--(d); \draw[ar] (d)--(e);
\end{tikzpicture}\\[2pt]
{\scriptsize\itshape Green: exact motif ([E]).\quad Gold: physical bridge ([C]) --- the same ratio, two epistemic types.}
\end{center}}

% ===== B7: the K5 action-lattice mnemonic ========================
%  Layout: complete graph K5 (pentagon, all 10 edges) on the LEFT, the
%  Pascal-row reading list on the RIGHT, vertically centred on the graph
%  so the two blocks never overlap.
\newcommand{\TFPTactionkfive}{%
\begin{center}
\begin{tikzpicture}[font=\scriptsize]
  % --- K5 vertices on a regular pentagon, centred at the origin ---
  \foreach \i in {1,...,5}{\coordinate (k\i) at ({90+72*(\i-1)}:12mm);}
  \begin{scope}[on background layer]
    \foreach \i in {1,...,5}{\foreach \j in {1,...,5}{\ifnum\i<\j
      \draw[tfgray!55,line width=0.5pt] (k\i)--(k\j);\fi}}
  \end{scope}
  \foreach \i in {1,...,5}{\fill[tfblue!70!black] (k\i) circle (1.8pt);}
  \node[tfgray,font=\scriptsize\itshape] at (0,-1.9) {$K_5$};
  % --- reading list, anchored to the right of the whole graph, centred ---
  \node[anchor=west,align=left,inner sep=0pt] at (1.9cm,0) {%
    $\tbinom51=5$\ \ vertices --- action $x^1$ (Hubble slot)\\[1pt]
    $\tbinom52=10$\ edges --- $\Lambda$ density pair $x^{10}$\\[1pt]
    $\tbinom53=10$\ triangles --- dual sector\\[1pt]
    $\tbinom54=5$\ \ complements --- inverse sector\\[1pt]
    $\tbinom55=1$\ \ determinant --- closure};
\end{tikzpicture}\\[2pt]
{\scriptsize\itshape $K_5$ on the five carrier slots; the action powers $x^{1},x^{5},x^{10}$ are its
$\tbinom5k$ faces \ \textcolor{tfgray}{(mnemonic; the physical scale map stays [C]).}}
\end{center}}


\newcommand{\cthree}{c_3}
\newcommand{\vgeo}{v_{\mathrm{geo}}}
\newcommand{\Mpl}{\bar M_{\mathrm{Pl}}}
\newcommand{\Nfam}{N_{\mathrm{fam}}}
\newcommand{\gcar}{g_{\mathrm{car}}}
\providecommand{\PP}{\mathbb{P}}
\providecommand{\Z}{\mathbb{Z}}
\setlist{itemsep=2pt,topsep=3pt}
\newcolumntype{Y}{>{\raggedright\arraybackslash}X}

\title{\vspace{-1.2cm}\textbf{TFPT --- Red Team / Stress Test layer}\\[2mm]
\large Five adversarial targets (A--E): try to \emph{break} the reductions, not confirm them}
\author{\TFPTauthors}\date{\TFPTdatestamp}

\def\TFPTdoctitle{Red Team}
\begin{document}
\maketitle

% ---------------------------------------------------------------------
%  tfpt_docset.tex -- THE canonical TFPT document-set box.
%  Single source of truth for the document list; \input by the
%  introduction and every paper. Uses only universally available
%  macros (keybox + plain LaTeX math) so it compiles in every doc.
%  Each document adds its own "You are reading ..." line directly
%  after the \input.  Numbering is aligned with the file names:
%  document N == tfpt_N_*, the introduction is the entry point, and
%  R/H/O are the appendix-style companions.
% ---------------------------------------------------------------------
\begin{navbox}[The TFPT document set --- what is treated where]
\textbf{Plain language:} TFPT (Topological Fixed-Point Theory) is a small discrete compiler. Two
inputs --- the seam constant $c_3=\tfrac1{8\pi}$ and the carrier rank $g_{\mathrm{car}}=5$ --- build
$D_5\oplus A_3+\mu_4\Rightarrow E_8$ and read off the Standard-Model skeleton, the constants and the
scale grammar. The series is \textbf{the introduction plus five numbered papers and three companions},
best read in order:
\begin{enumerate}[leftmargin=2.0em,itemsep=1pt,label=\textbf{\arabic*.}]
\setcounter{enumi}{-1}
\item \texttt{introduction} --- reading guide, compiler closure, predictions, the dependency DAG and the
      single proof ledger (\texttt{verification/status\_ledger.csv}). \emph{(entry point)}
\item \texttt{tfpt\_1\_architecture\_e8} --- the two axioms, the derivation map, the EM fixed point
      $\alpha^{-1}$, the $D_5\oplus A_3+\mu_4\Rightarrow E_8$ construction, the QBL theorem chain.
\item \texttt{tfpt\_2\_standard\_model} --- the SM in one $\varphi_0$-ladder formula, flavor from
      parabolic transport (sheet diamond, dual anchor, $Q_+$ cohomology), the worked closures.
\item \texttt{tfpt\_3\_e8\_audit\_bootstrap} --- the seven $E_8$ slices as an audit raster, the cascade
      bridge, the M\"obius bootstrap.
\item \texttt{tfpt\_4\_frontier} --- honest status of $\eta_B$, $m_p/m_e$, Koide, dark matter, full QG.
\item \texttt{tfpt\_5\_redteam} --- the adversarial audit: declared attacks, what survives, what each
      reduces to, and the kill tests.
\end{enumerate}
\noindent\textbf{Companions:} \textbf{R.} \texttt{tfpt\_research\_contracts} --- the open interfaces
$(v_{\mathrm{geo}}, G_{\mathrm{net}}, F_{\mathrm{transfer}})$ as numbered contracts;
\textbf{H.} \texttt{tfpt\_horizon\_readouts} --- Appendix~H (unit-system reframe): one seam constant as
the universal horizon thermal code, the Nariai/anchor surface, the resummed seam clock;
\textbf{O.} \texttt{origin\_theory} --- why no free fundamental number remains (exact core $+$ one
honestly-typed cyclic interpretation).
\end{navbox}

\noindent\textbf{You are reading \textbf{document~5} (Red Team).}

% ---------------------------------------------------------------------
%  tfpt_status.tex -- THE canonical "status at a glance" card.
%  Single source for the overall theory status; \input near the front
%  of every document so there is ONE status statement, not one per
%  paper.  Requires tfpt_style.tex (keybox, markers).  Detailed,
%  paper-specific status expansions live only where they belong
%  (introduction: five-questions + closure ledger; tfpt_4: frontier
%  table; tfpt_research_contracts: the contracts).
% ---------------------------------------------------------------------
\begin{keyboxnb}[TFPT status at a glance --- read this card the same way in every document]
\textbf{Two inputs, one compiler.} From the boundary constant $c_3=\tfrac1{8\pi}$ (axiom~P1) and the
five-slot carrier $g_{\mathrm{car}}=5$ (axiom~P2) --- jointly the elementary symmetric data of one anchor
$a=(1,1,2)$ --- the discrete compiler $D_5\oplus A_3+\mu_4\Rightarrow E_8$ is \emph{built}, and the
Standard-Model skeleton (gauge group, three families, hypercharges, the flavor matrix), the dimensionless
constants ($\alpha^{-1}=137.0359992$) and the scale grammar are read off as fixed points. The algebraic
core is machine-checked; physical readouts run through \emph{named bridges} whose status is typed below.
\smallskip\\
\noindent\textbf{What is closed, and what genuinely remains.} The discrete/algebraic layer is closed
(\mE); the everything-open residual reduces to \emph{three named interface problems}:
\begin{center}\small
\renewcommand{\arraystretch}{1.18}
\begin{tabular}{@{}lll@{}}
\toprule
\textbf{Interface} & \textbf{Question} & \textbf{Status} \\
\midrule
$v_{\mathrm{geo}}$ & the one metrology unit ($=1/\sqrt G=m/\mu$); No-Unit Thm: no compiler scale & primitive \mO \\
$G_{\mathrm{net}}$ & simple-current extension $(D_5)_1{\otimes}(A_3)_1{\rtimes}\langle(1,1)\rangle\cong(E_8)_1$ & \mE/\mC \\
$F_{\mathrm{transfer}}$ & one functor, four typed interfaces (Koide, $\eta_B$, axion, $m_p/m_e$) & \mC \\
\bottomrule
\end{tabular}
\end{center}
\noindent Everything carries a claim ID resolving to a row of \texttt{verification/status\_ledger.csv}
(the single source of truth); \textbf{if the text and the ledger ever disagree, the ledger wins.} The
historical labels $(U_{\mathrm{wall}},G_{\mathrm{metric}},F_{\mathrm{frontier}})$ are retained only for
ledger continuity --- the live residual is $v_{\mathrm{geo}}\oplus G_{\mathrm{net}}\oplus F_{\mathrm{transfer}}$.
\markerkey
\end{keyboxnb}

\TFPTbadges
\TFPTclaimstack[none]

\vspace{-0.6cm}

\begin{keybox}[What this note is]
This is the deliberately \emph{adversarial} layer requested in the \vref{v78} review. Its purpose is
\textbf{not} to confirm TFPT but to attack the five load-bearing reductions at their weakest
logical transitions. Each target is run through one fixed protocol --- minimal statement,
assumptions, logical chain, validity conditions, counterexample search, limiting/degenerate
cases, alternative structures, provisional verdict --- and produces a deliberately narrow output:
where the statement works, where it could fail, and the residual risk. The layer is machine-backed
by \veri{redteam/rt\_A\_e8net.py} \dots \veri{redteam/rt\_E\_vgeo.py}, \veri{redteam/rt\_F\_qft4d.py} and
\veri{redteam/run\_redteam.py}; a red-team check asserts an \emph{adversarial} fact (a counterexample
really exists, a hidden assumption is really needed, a firewall really holds), and the honest
outcome lives in the \textbf{status} of each target, never in a green pass.
\end{keybox}

\begin{center}\small
\renewcommand{\arraystretch}{1.3}
\begin{tabularx}{\textwidth}{@{}l Y l Y@{}}
\toprule
\textbf{Target} & \textbf{Minimal claim} & \textbf{Status} & \textbf{Residual risk} \\
\midrule
A & seam--Calder\'on measure $=$ $(E8)_1$ net & \textcolor{tfred}{\textbf{reduced, not closed}} & \emph{one} statement: boundary-net holomorphy $+$ $c{=}8$ ($\Leftrightarrow$ the index-$4$ inclusion); $E_8$ and bulk uniqueness then automatic (\vref{v83}/\vref{v87}/\vref{v89}; Lie-level realisation \vref{v143}). Since reduced to \emph{zero new} open content: the free-bulk premise is a fixed-point theorem, the cone is eliminated (cone gap $=$ one-particle gap), and the residual factors into the already-open $A_2$ (net assembly) $+$ \texttt{GATE.QGEO} (\vref{v160}--\vref{v165}). The three-residual form below is the historical development. \\
B & $\gcar=5$ forced by the Pascal rule & \textcolor{tfgold}{\textbf{survives (narrowed)}} & boundary derivation of half-spinor exhaustion (degree-2 truncation) \\
C & $k=\cthree/2=\tfrac{1}{16\pi}$, $S/A=\tfrac14$ & \textcolor{tfgreen}{\textbf{survives}} & absolute $1/G$ (UV-sensitive) is the one anchor \\
D & ratios$+$products $\Rightarrow$ one scale $\vgeo$ & \textcolor{tfgold}{\textbf{survives (narrowed)}} & CP phases ($\delta_{\mathrm{CKM/PMNS}}$) not covered by $\vgeo$ \\
E & one scale $\vgeo$ carries the theory & \textcolor{tfgold}{\textbf{survives (narrowed)}} & EW/reheating/leptogenesis scales $+$ seam$=$Planck identification \\
\bottomrule
\end{tabularx}
\end{center}

% =====================================================================
\section*{Method}
% =====================================================================
The red team treats each reduction as a hostile witness. A confirmatory script that always passes
is worthless here: the layer is built so that it \emph{may} downgrade a claim. Three outcomes are
allowed --- \textbf{survives} (stands as worded), \textbf{survives (narrowed)} (stands only after a
silent assumption is made explicit), and \textbf{reduced, not closed} (the conservative wording is
the correct one). A fourth, \textbf{broken}, is reserved for an actual failure; none occurred.

% =====================================================================
\section*{Target A --- the $(E8)_1$ boundary-net identification}\label{sec:A}
% =====================================================================
\noindent\textbf{Minimal claim.} The ambient metric/projective measure is reduced to identifying the
seam--Calder\'on boundary measure with the $(E8)_1$ lattice net, carrier subnet
$(D5)_1\times(A3)_1$.\veri{redteam/rt\_A\_e8net.py}

\smallskip\noindent\textbf{Logical chain (established).} The Sugawara level-1 central charges are
$c(E8)_1=\tfrac{248}{31}=8$, $c(D5)_1=5=\gcar$, $c(A3)_1=3=\Nfam$, so the conformal-embedding
criterion $c_{\mathrm{coset}}=c(E8)-c(D5)-c(A3)=0$ holds \mE. This is \emph{compatibility}.

\begin{gapbox}[Where it breaks: central charge underdetermines the net]
Counterexample: $(D8)_1=SO(16)_1$ also has $c=\tfrac{120}{15}=8$, but is \emph{not} holomorphic
(four primaries $1,v,s,c$) --- a distinct $c=8$ net. Hence equal central charge is \emph{necessary,
not sufficient}: $(E8)_1$ is the unique \emph{holomorphic} $c=8$ net (even unimodular rank-8
lattice $=E8$), so \textbf{holomorphy (single vacuum sector, $\mu$-index $1$) is the load-bearing
extra assumption}. Dropping chirality is worse: the $c=8$ Narain family is positive-dimensional.
The constructive map seam--Calder\'on kernel $\to$ $(E8)_1$ and the uniqueness of bulk
reconstruction are not established; the gap $\Delta_{\mathrm{eff}}>0$ \emph{supports} tightness
(clustering) but does not prove it.
\end{gapbox}

\noindent\textbf{Verdict.} Keep ``reduced to a rigorous boundary-net identification problem''; do
\emph{not} write ``metric sector closed''. The red team confirms the conservative wording. Status
\mC. Residual: holomorphy proof $+$ constructive boundary map $+$ bulk-reconstruction uniqueness.

% =====================================================================
\section*{Target B --- carrier rank / the Pascal condition}\label{sec:B}
% =====================================================================
\noindent\textbf{Minimal formal claim.} $2^{g}=g^2+g+2$ has the unique solution $g=5$ (Lean-verified
\mE). This arithmetic is \emph{not} attacked. The attack is upstream: why must the carrier obey the
Pascal half-spinor exhaustion $2^{g-1}=\binom{g}{0}+\binom{g}{1}+\binom{g}{2}$?\veri{redteam/rt\_B\_pascal.py}

\begin{keybox}[Where it narrows: the rule, not the arithmetic, selects 5]
Perturbing the constant breaks it: $2^{g}=g^2+g+c$ gives $\{5\}$ only at $c=2$; neighbours give
other or empty solution sets. The truncation \emph{degree} is the real choice: with
$2^{g-1}=\sum_{k\le K}\binom{g}{k}$ one finds $K{=}0\!\to\!g{=}1$, $K{=}1\!\to\!3$,
$K{=}2\!\to\!5$, $K{=}3\!\to\!7$, i.e.\ $g=2K+1$. So choosing $K=2$ (to reach $\gcar=5$) is a
physical postulate. It is corroborated --- the $D5$ half-spinor $\dim S^+=2^{g-1}=16\Leftrightarrow g=5$
and the family closure $\Nfam=(2^{g-1}{-}1)/g$, $g+\Nfam=8$ both pick $5$ --- but each is itself a
postulate, not a theorem.
\end{keybox}

\noindent\textbf{Verdict.} Theorem A's arithmetic core stands; the Pascal \emph{selection} must be
typed \mO/\mC, not \mE. Residual: a boundary/seam derivation of half-spinor exhaustion.

% =====================================================================
\section*{Target C --- $k=\cthree/2$ and the seam-area coefficient}\label{sec:C}
% =====================================================================
\noindent\textbf{Minimal claim.} In reduced units $k=\cthree/2=\tfrac{1}{16\pi}$ and
Fursaev--Solodukhin gives $S/A=\tfrac14$.\veri{redteam/rt\_C\_kc3.py}

\begin{winbox}[The dimensional firewall holds]
The only legal chain is
\[
  k_{\mathrm{red}}=\tfrac{\cthree}{2}=\tfrac{1}{16\pi}\ \text{(dimensionless)},\qquad
  k_{\mathrm{phys}}=\tfrac{k_{\mathrm{red}}}{G},\qquad
  S=4\pi\,k_{\mathrm{phys}}A_{\mathrm{phys}}=\tfrac{A_{\mathrm{phys}}}{4G}.
\]
The forbidden form $k_{\mathrm{phys}}=\cthree/2$ is dimensionally false (legal only at $G=1$). A
scan of \texttt{v67/\vref{v68}/v73} finds every $\cthree/2$ used as the \emph{reduced} coefficient and
\textbf{zero} naked dimensionful identities; the firewall's teeth are verified against a synthetic
violation. Internally consistent with $1/G$ being UV-sensitive (Sakharov/Connes, \vref{v68}).
\end{winbox}

\noindent\textbf{Verdict.} Safe as worded (reduced coefficient). Status \mE\,(structure) $+$ \mO\,for
the residual. Residual: the absolute $4D$ Newton scale $1/G$ ($\propto \Lambda^2 f_2$) stays the one
dimensionful anchor; nothing \emph{dimensionless} is open.

% =====================================================================
\section*{Target D --- $U_{\mathrm{point}}\to\vgeo$}\label{sec:D}
% =====================================================================
\noindent\textbf{Minimal claim.} Ratios plus sector products determine the dimensionless amplitudes,
leaving one common scale $\vgeo$.\veri{redteam/rt\_D\_upoint.py}

\begin{keybox}[Where it narrows: the bijection needs explicit hypotheses]
For positive reals, (ratios, product) $\Leftrightarrow$ (individuals) is a bijection, and the TFPT
sectors satisfy it (positive rational $c$, odd size $3=\Nfam$). But it \emph{fails} off-assumption:
even sector size leaves a sign ambiguity ($\{2,3\}$ vs $\{-2,-3\}$ share ratio $3/2$ and product
$6$); complex amplitudes leave an $n$-th-root branch ambiguity ($\{1,1,1\},\{\omega,\omega,\omega\},
\{\omega^2,\omega^2,\omega^2\}$ share all ratios and product $1$); a zero amplitude is degenerate.
\textbf{Genuine residual:} the bijection fixes amplitude \emph{magnitudes} only --- CP phases
$\delta_{\mathrm{CKM}},\delta_{\mathrm{PMNS}}$ are complex and are \emph{not} folded into $\vgeo$.
\end{keybox}

\noindent\textbf{Verdict.} Holds for TFPT once four hypotheses (real, positive, fixed branch, no
zeros / odd sector size) are made explicit; they are silent in \vref{v75} and must be named. Status \mE\,
(reduction) $+$ \mO\,($\vgeo$). Residual: the CP-phase sector.

% =====================================================================
\section*{Target E --- $\vgeo$ as the unique dimensional floor}\label{sec:E}
% =====================================================================
\noindent\textbf{Minimal claim.} No pure-number theory can derive an absolute dimensionful scale; all
scales are ratios to one unit $\vgeo$.\veri{redteam/rt\_E\_vgeo.py}

\begin{keybox}[Single-scale audit: exact for the certified tiers, conditional beyond]
The closed compiler $+$ protected IR readouts all reduce to (pure number)$\times\vgeo$:
$M_{\mathrm{scal}}=\cthree^{7/2}\Mpl$, $f_a=(\cthree/|\mu_4|)^{7/2}\Mpl$, masses
$=(\pi/\sqrt2)\,c_f\,\varphi_0^{k}\,\vgeo$, and $1/G$ shares the same anchor (\vref{v68}/\vref{v75}). The audit
\emph{surfaces} candidates that are not pure reductions: the seam cutoff $\Lambda$ (only an
\emph{identification} seam$=$Planck collapses it to $\vgeo$), the electroweak matching scale, the
reheating temperature, and the leptogenesis scale --- all frontier \mC/\mO\ inputs that must not be
silently absorbed into $\vgeo$.
\end{keybox}

\noindent\textbf{Verdict.} $\vgeo$ is the unique dimensional floor of the certified tiers (true by
dimensional analysis). Full single-scale-ness is conditional on the flagged identifications staying
honestly typed. Status \mE\,(ratios) $+$ \mO\,(one scale). Residual: explicit reduction (or honest
frontier typing) of the EW / reheating / leptogenesis scales $+$ the seam$=$Planck cutoff.

% =====================================================================
\section*{Target F --- the perturbative 4D-QFT $+$ scale round (v269--v275)}\label{sec:F}
% =====================================================================
\noindent\textbf{Minimal claim.} The round published in release 5.2 --- the Epstein--Glaser
$S_{\mathrm{pert}}$ layer (\vref{v269}/\vref{v271}/\vref{v273}), the PMNS Jarlskog CP strength
(\vref{v270}), the absolute $\nu$-scale typing (\vref{v272}), the scale over-determination
(\vref{v274}) and the \texttt{QG.AMB.01} roadmap (\vref{v275}).\veri{redteam/rt\_F\_qft4d.py}

\begin{keybox}[Two attacks land; the round survives once they are named]
\textbf{(1) The gravity $S_{\mathrm{pert}}$ is non-unitary.} Power-counting renormalizability is \emph{not}
unitarity: the $R^2/$Weyl$^2$ sector gives a four-derivative propagator $1/(p^2(p^2+M^2))$ with a
negative-residue massive spin-2 mode (the Stelle ghost). So $S_{\mathrm{pert}}$ is unitary for the
matter$+$gauge sector only; the gravity perturbative S-matrix carries the ghost --- which is exactly why the
\emph{nonperturbative} \texttt{QG.AMB.01} is the real frontier, not a new gap.
\textbf{(2) The over-determination is conditional.} \vref{v274}'s Route-2 inverts the
$\rho_\Lambda/\Mpl^4=(3/4\pi^2)e^{-2\alpha^{-1}}$ prediction, which is itself \texttt{LAMBDA.BRANCH.01} \mC;
so the $0.11\%$ gravity-vs-cosmology agreement is conditional on the $\Lambda$-branch (the gravity route is
unconditional). \textbf{Already-honest:} $J_{\mathrm{PMNS}}$ inherits $\delta$'s \mC\ provenance (assembled
$\mu_6$ phase, not an $M_\nu$ output) and the $\nu$-scale is \mO\ with the NO floor $\Sigma m_\nu=0.0586$ eV as
the one kill test.
\end{keybox}

\noindent\textbf{Verdict.} \textsc{Survives (narrowed)}: the round stands as typed; the two narrowings above are
now folded into \vref{v269} and \vref{v274}. Residual: the perturbative gravity S-matrix is non-unitary (the
Stelle ghost) --- this \emph{is} the \texttt{QG.AMB.01} frontier; and the anchor over-determination is
conditional on the $\Lambda$-branch.

\par\smallskip\noindent\textbf{Under examination: an infinite-derivative narrowing (\mC, does \emph{not}
close the gate).} The Stelle ghost is a feature of the \emph{local} $R{+}R^2/$Weyl$^2$ \emph{truncation}:
the four-derivative propagator $1/(p^2(p^2{+}M^2))=(1/M^2)[1/p^2-1/(p^2{+}M^2)]$ carries the
negative-residue spin-2 mode. But that truncation is the $a_4$ Seeley--DeWitt order of the spectral action
$\mathrm{Tr}\,f(D^2/\Lambda^2)$, whose cutoff is \emph{not} a free scheme --- it is the $\beta{=}1$ seam KMS
weight $f(u)=e^{-u}$ (\vref{v259}, no new parameter). Keeping that exponential cutoff \emph{un-truncated}
points to a nonlocal, infinite-derivative graviton form factor $p^2 a(p^2)$; for an \emph{entire} factor
$a(p^2)=e^{p^2/M^2}$ (nowhere zero) the dressed propagator $1/(p^2 a)$ keeps its \emph{only} pole at $p^2{=}0$
--- no ghost (Tomboulis 1997; Biswas--Mazumdar--Siegel 2006) --- whereas any \emph{polynomial}
(finite-derivative) truncation has a real zero $=$ the ghost pole. So \emph{if} the resummed form factor is
entire, the perturbative ghost is a truncation artefact: a conditional \mC\ narrowing of this attack. It does
\emph{not} close \texttt{QG.AMB.01} --- the trace regulator $f$ is not yet shown to resum to an entire kinetic
form factor $a(\Box)$, and a ghost-free \emph{perturbative} graviton is still not the nonperturbative ambient
measure. (\texttt{QGAMB.IDG.01}, \mC/\mO; consistent with the \texttt{SEAM.EQUIV.01} firewall.)
\veri{v304\_idg\_nonlocal\_ghost.py}

% =====================================================================
\section*{Outcome}
% =====================================================================
\par\smallskip\noindent\textbf{What the red team changed.}
No target broke. The hardest gate (A) is confirmed \emph{open} with the conservative wording;
four targets (B, D, E, F) \emph{survive once a silent assumption is named}; one (C) survives as
worded. The net effect is not new physics but sharper typing: the package now states \emph{where}
each reduction would fail and \emph{which} assumptions are truly necessary, exactly at the weakest
logical transitions. Reproduce with \texttt{cd verification/redteam \&\& python run\_redteam.py}.

\begin{winbox}[Follow-up --- the verdicts were then used as a worklist \mE\ (\veri{v83\_e8net\_holomorphic\_uniqueness.py})]
Two residuals were attacked directly after the red-team pass; the result is recorded in ledger rows
\texttt{GATE.METRIC.04} and \texttt{CAR.PASCAL.01}.

\smallskip\noindent\textbf{Target A, residual 1 --- closed \mE.} At level 1 the number of primaries equals
$\det(\text{Cartan})=|\text{fundamental group}|$: $A_3{=}4$, $D_5{=}4$, $D_8{=}4$, $\mathbf{E_8{=}1}$. So holomorphy
(single vacuum sector, $\mu$-index $1$) is \emph{necessary} --- it excludes the same-$c$ rival $(D_8)_1=SO(16)_1$
($c=120/15=8$ but four primaries $1,v,s,c$) --- \emph{and sufficient}: a holomorphic $c=8$ chiral CFT is the lattice
theory of an even unimodular rank-$8$ lattice, of which there is exactly one, $E_8$, by the Minkowski--Siegel mass
formula ($\text{mass}=1/|W(E_8)|=1/696729600$, and $|\mathrm{Aut}(E_8)|$ saturates it). The
$(D_5)_1\times(A_3)_1$ embedding $c=5+3=8$ stays \emph{compatibility} ($16$ primaries vs $1$). \textbf{Consequence:}
the constructive map need only show the seam--Calder\'on boundary net is holomorphic with $c=8$ (then $E_8$ is
automatic), so Target~A drops from \emph{three} residuals to \emph{two}: (i)~boundary-net holomorphy $+\,c=8$
($\Delta_{\mathrm{eff}}>0$ supports, not proves), and (ii)~bulk-reconstruction uniqueness --- both still \mC/\mO.

\smallskip\noindent\textbf{Target B --- reduced \mE.} The ``$K=2$ Pascal truncation'' is not free:
$K=(g-1)/2$ is the Pascal-row midpoint $\sum_{k\le(g-1)/2}\binom{g}{k}=2^{g-1}$ (odd $g$), i.e.\ the half-spinor
split; the carrier is the even Clifford half-spinor $\Lambda^{\mathrm{even}}(\mathbb C^{5})=1+10+5=16=\dim S^+$. So
B's residual reduces to the single standard input ``carrier $=$ half-spinor of $\mathrm{Spin}(10)$'' (the Lean
arithmetic core stays \mE).
\smallskip\\
\noindent\textbf{Target B --- a backward certificate (\veri{v219\_icosahedral\_mckay.py}).} The choice
$(|\Z_2|,\Nfam,\gcar)=(2,3,5)$ that B attacks is, \emph{given} the $E_8$ closure, the icosahedron: by the
McKay correspondence $E_8$ is the exceptional top of the finite-$SU(2)$ tower ($2T{\to}\hat E_6$,
$2O{\to}\hat E_7$, $2I{\to}\hat E_8$), the binary icosahedral group $2I$ (order $120=|R^+(E_8)|$) has irrep
degrees $\{1,2,2,3,3,4,4,5,6\}$ equal to the affine-$E_8$ Kac marks ($\sum=30=h(E_8)$), and it contains the
$2,3,5$ rotation-axis orders. This does \emph{not} break B's verdict --- it is a \emph{backward} certificate
(it presupposes $E_8$), so $\gcar=5$ stays the P2 interface \mO; what it removes is the impression that
$(2,3,5)$ is an arbitrary triple, since no spherical McKay node sits above the icosahedron.

\smallskip\noindent Still open, typed honestly: A\,(i)$+$(ii); Target~D's CP-phase residual (at $\theta_{13}=0$ the
Dirac phase is undefined, so this needs a texture correction); Target~E's reheating/leptogenesis scales and the
seam$=$Planck identification. Target~C is a dimensional \emph{floor}, not a gap. \textbf{No target was promoted.}
\end{winbox}

\begin{winbox}[Second follow-up round \mE\ (\veri{v87\_bulk\_uniqueness\_reduction.py}, \veri{v88\_cp\_phase\_audit.py}, \veri{v86\_nstar\_reheating.py})]
\smallskip\noindent\textbf{Target A, residual (ii) --- conditionally closed \mE: Target A $=$ ONE residual.}
Bulk-reconstruction uniqueness is \emph{not} independent of residual (i). By the algebraic-QFT classification of
full 2D CFTs over a chiral net (Longo--Rehren; Kawahigashi--Longo--M\"uger; Bischoff--Kawahigashi--Longo--Rehren),
the possible bulks are the haploid commutative Q-systems $\simeq$ physical modular invariants of $\mathrm{Rep}(A)$.
For a \emph{holomorphic} net $\mathrm{Rep}(A)=\mathrm{Vect}$, so the bulk pairing is \emph{unique}. The contrast is
machine-checked: the same-$c$ rival $SO(16)_1$ (discriminant category $\mathbb Z_2{\times}\mathbb Z_2$, $\mu$-index $4$) admits
\emph{six} modular invariants --- diagonal, $s{\leftrightarrow}c$ swap, both $E_8$-extension invariants
$|\chi_0{+}\chi_{s/c}|^2$ (the lattice glue $E_8=D_8\cup(D_8{+}s)$, $h_s{=}1$ bosonic), and two twisted pairings ---
so without holomorphy the $c=8$ bulk is \emph{multiply} ambiguous; with it, trivially unique. Hence the entire
Target~A reduces to the single theorem ``seam--Calder\'on boundary net is holomorphic with $c=8$'' (\mC/\mO;
ledger \texttt{GATE.METRIC.05}). The theorem also has an equivalent, more tractable \emph{index} form
(\veri{v89\_carrier\_index\_lemma.py}, ledger \texttt{GATE.METRIC.06}): by Kawahigashi--Longo--M\"uger,
$\mu_A=[B{:}A]^2\mu_B$, the carrier inclusion has Jones index
$[(E_8)_1:(D_5)_1{\times}(A_3)_1]=\sqrt{16}=4=|\mu_4|$ --- the glue-group order \emph{is} the inclusion
index --- and the extension is the $\mu_4$ simple-current extension (all three glue sectors are $h=1$
currents; $248=45{+}15{+}64{+}64{+}60$). Since an index-$4$ simple-current extension of the carrier has
$\mu=16/4^2=1$, \textbf{holomorphy follows from the index statement}: Target~A $\Leftrightarrow$ ``the seam
net contains the carrier net with Jones index $4$ as its $\mu_4$ simple-current extension'' --- both ends
constructed objects, the index controllable by Longo theory, the gap supporting finiteness. And \emph{which}
extension carries no freedom (\veri{v92\_glue\_uniqueness.py}, ledger \texttt{GATE.METRIC.07}): exhaustive
classification of the carrier discriminant form $(\mathbb Z_4{\times}\mathbb Z_4,\,q{=}(5x^2{+}3y^2)/8)$ shows the
extension tower is \emph{completely rigid} --- exactly two Lagrangian glues (the two chiralities, identified
by the sheet $\mathbb Z_2$; the same two objects as the $E_8$-extension invariants above), and exactly one halfway
extension, whose induced form $q=\{0,0,0,\tfrac12\}$ \emph{is} the $D_8$ discriminant form: the $SO(16)_1$
rival is the unique intermediate of the tower, not an arbitrary counter-model. Tower:
carrier\,$(\mu{=}16)\to SO(16)_1\,(\mu{=}4)\to(E_8)_1\,(\mu{=}1$, sheet pair$)$ --- nothing else exists.
Target~A is thereby the \emph{bare} index statement.

\smallskip\noindent\textbf{Target A, final form --- the Simple-Current Extension Theorem \mE/\mC\
(\veri{v154\_simple\_current\_theorem.py}).} Assembled into one statement: the carrier net
$A=(D_5)_1{\otimes}(A_3)_1$ extended by the isotropic glue $L=\langle(1,1)\rangle$ has $[B:A]{=}|L|{=}4$,
$c{=}5{+}3{=}8$, $\mu(B){=}\mu(A)/|L|^2{=}16/16{=}1\Rightarrow$ holomorphic $\Rightarrow B\cong(E_8)_1$. The
\emph{entire} adversarial residual of Target~A is now the single seam-side identification ``the RP
seam--Calder\'on net \emph{is} $A$'' --- the algebra ($B\cong(E_8)_1$) is exact and machine-checked
(\vref{v89}/\vref{v92}/\vref{v125}/\vref{v143}/\vref{v148}). That residual is now sharply located: the
analytic core (``the kernel is quasi-free'') \emph{follows} from a free RP bulk --- the boundary marginal of
a Gaussian bulk is the compression of a contraction (\vref{v155}), the \emph{same} premise the area law and
the EH mechanism use --- and the net is \emph{explicitly constructed and verified}: DtN $=|k|$, $16$ free
Majoranas, currents $248{=}120{+}128$, holomorphic character $E_4/\eta^8=j^{1/3}$ with $248$ at level $1$
(\vref{v156}). \emph{Final reduction (\vref{v160}--\vref{v165}):} the free-bulk premise is recast as a
fixed-point theorem (quasi-free $\Rightarrow$ all higher cumulants $\kappa_{2n}{=}0$, \vref{v160}); the
infinite Schwinger cone is \emph{eliminated} --- by Bogoliubov second quantisation the cone gap equals the
one-particle gap $(2/3)^6$ (\vref{v161}), whose value is forced by $\mu_4$-equivariance (\vref{v162}); and
the remaining $16$-dim identification factors into the two \emph{already-open} items --- $A_2$ net existence
(an assembly of established net constructions, \mC) and $R_1=$ \texttt{GATE.QGEO} (``seam boundary $=
\mathbb P^1{\setminus}\mu_4$'') --- with the irreducible core $\{\pi,v_{\mathrm{geo}}\}$ a theorem
(\vref{v165}). So Target~A carries \emph{no new independent open content}: it is a unification, not an
unconditional proof. \emph{The two structural residuals are then discharged to a theorem
(\vref{v175}):} the CAR second-quantisation functor $\Gamma(t){=}\bigoplus_m\Lambda^m(t)$ makes full-cone
reflection positivity reduce \emph{for every $m$} to the one-particle contraction (verified on the complete
$2^{16}$-dim Fock space: contraction, fixed multiplicity $2^8$, sub-leading $=(2/3)^6$), and $A_2$ is an
assembled and verified $(E_8)_1$ certificate (character $1{+}248q{+}4124q^2{+}\dots$, embedding
$248{=}120{+}128$, $E_8$ Cartan even unimodular). Both then need the \emph{same} single input --- the seam
realisation \texttt{QGEO.REALIZE.01} --- which stays the one irreducible \mO; net existence and full-cone RP
become \mE. \emph{That single residual is then driven to bedrock (\vref{v176}--\vref{v181}):} it is
assembled as one central theorem (\vref{v176}), split into the marks $+$ kernel obligations whose finite
cores close (\vref{v177}/\vref{v178}), unified into one conformal-realisation premise (\vref{v179}), and
reduced --- via uniformisation, Ker\'ekj\'art\'o and the order-$4$ M\"obius classification --- to a seam-%
\emph{isometry} premise (\vref{v180}) and finally to the strictly weaker conformal-\emph{symmetry} / deck
premise \texttt{QGEO.SYM.01} (\vref{v181}): the carrier $\mu_4$ clock is the conformal deck of the seam.
Below that the residual is \emph{definitional}, not a missing theorem --- the single irreducible
structural residual of the whole theory. It is since driven to its sharpest, \emph{non-circular} form
(\vref{v194}--\vref{v198}): the raw seam DtN is RP-canonical (Osterwalder--Schrader on the quasi-free
state), and the bedrock is cracked to the \emph{state-invariance} $\omega\circ\rho=\omega$ --- the DtN
principal symbol commutes \emph{exactly} with the clock, and Tomita--Takesaki gives
$[\rho,\Lambda_\Sigma]=0$ with no conformal-covariance assumption, removing the Bisognano--Wichmann
circularity --- the maximally-operational form of the same postulate, still \mO. \emph{The geometric and
conditional cores are now Lean-formalised} (\textsc{audit: pass}): the \textsc{Uniform} normal form
($z\mapsto iz$, $\sigma\rho\sigma=\rho^{-1}$, orbit$\to\mu_4$; \texttt{FORM.QGEO.02}) and the implication
mark-local$\Rightarrow\omega\circ\rho=\omega$ (\texttt{FORM.QGEO.01}) --- so the deduction below the
premise is machine-proved, while \texttt{QGEO.SYM.01} itself stays the one \mO\ postulate.

\smallskip\noindent\textbf{Target D --- the CP-phase residual quantified \mE\ (not closed).} With the frozen leading
reading $\delta=\pi/3+3\lambda_C^2=68.65^\circ$: pull vs $\gamma_{\mathrm{PDG}}=65.7^\circ\pm3.0^\circ$ is $+0.98\sigma$
(survives). Audit, \emph{not} promoted: the data central value sits $0.07^\circ$ from the alternative coefficient
reading $\pi/3+2\lambda_C^2$ ($|\mathbb Z_2|$ instead of $\Nfam$) --- a textbook look-elsewhere trap; the registry keeps
coefficient $3$ (\texttt{REG.FREEZE.01}). Decision threshold: the readings differ by $\lambda_C^2=2.88^\circ$, i.e.\
$3\sigma$-distinguishable once $\sigma_\gamma\le0.96^\circ$. The $J$-inversion is flagged magnitude-contaminated
(source-scale $s_{23}$ vs $M_Z$). Ledger \texttt{FLAV.CP.01}.

\smallskip\noindent\textbf{Target E --- one flagged scale computed \mC.} The reheating temperature is no longer a
silent input: $M_{\mathrm{scal}}=\cthree^{7/2}\bar M_{\mathrm{Pl}}$ \mE\ $+$ standard physics give
$T_{\mathrm{reh}}=9.6\times10^9$ GeV and $N_\star(0.05/\mathrm{Mpc})=51.4$ \mC\ (ledger \texttt{COSMO.NSTAR.01});
honestly recorded: this slow Higgs-channel point is \emph{$A_s$-disfavoured} ($-11.4\sigma$; the measured $A_s$
requires near-instantaneous reheating), so the point is conditional on the decay channel and the frozen band
$[50,60]$ stays the surface of record; the leptogenesis scale and the seam$=$Planck identification remain typed
inputs. \textbf{Again: no target was
promoted; one residual was \emph{merged} (A), one was \emph{quantified} (D), one input was \emph{computed} (E).}
\end{winbox}

\begin{winbox}[Target D --- the CP residual geometrically \emph{located} (not closed) \mE/\mC]
The follow-up round only \emph{quantified} the CP residual; this round \emph{locates} it geometrically,
in three consistent readings, without promoting it. \textbf{(i) A channel typing
(\veri{v227\_degree\_exponent\_channel\_split.py}).} The $E_8$ split $248=120+128$ is a
magnitude/phase split: the exponents sum to $120=|R^+(E_8)|$ (the \emph{magnitude} channel, where the
$(\text{ratios},\text{product})\Rightarrow\vgeo$ bijection lives), the fundamental degrees sum to
$128=\operatorname{rank}E_8\cdot\dim S^+$ (the \emph{phase/glue} channel). The un-covered CP phases of
Target~D are exactly the $128$ phase channel --- a clean home, not a new gate; whether $128$ also hosts a
dark sector is Frontier, \emph{not} asserted (this replaces the over-strong ``$S^-$ is dark matter''
reading). \textbf{(ii) A hexagonal modulus (\veri{v220\_cp\_hexagonal\_modulus.py}).} The seam
deck is the \emph{square} modulus ($j{=}1728$, $\Z/4$ clock; the kill-test selects it), which is real and
carries no phase; its exceptional partner is the \emph{hexagonal} modulus $j{=}0$ ($\tau{=}\rho{=}e^{i\pi/3}$,
$\Z/6$, CM by $\Z[\omega]$). The frozen leading reading $\delta=\pi/3+3\lambda_C^2=68.65^\circ$ has leading
term $\pi/3=\arg\rho$, the hexagonal CM phase; CP lives in the $\Z/6$ phase \emph{fiber} over the $\Z/4$ seam,
not in the deck. \textbf{(iii) An orientation (\veri{v225\_dual\_normal\_frame.py}).} The dual
normal frame $d=a^\top R^{-1}=-\tfrac12(1,1,-2)$ (Nariai/horizon normal) and $n=(5,-9,6)$ (first-generation
selector, $n^\top R=(8,0,0)$) has oriented volume $\det(\mathbf1,d,n)=21=\Nfam\cdot7$, and rotating by the
hexagonal unit gives $\operatorname{Im}\det(\mathbf1,d,\rho n)=21\sin(\pi/3)\neq0$ --- CP sits precisely in
the \emph{orientation} of the normal frame, exactly where the magnitude bijection fails. All three are \mE\
geometry; the physical CP \emph{derivation} stays the Target-D \mC\ residual (the seam deck is not
re-selected, the registry $\delta$ is untouched). \textbf{(iv) One hexagonal unit, two sheets
(\veri{v231\_cp\_mu6\_phases.py}).} The reduction is sharper still: \emph{both} CP phases are $\mu_6$ powers
of the single hexagonal CM unit $\rho=e^{i\pi/3}$ --- $\delta_{\mathrm{CKM}}^{\mathrm{lead}}=\arg(\rho^1)
=\pi/3$ (quark) and $\delta_{\mathrm{PMNS}}=\arg(\rho^4)=4\pi/3$ (lepton, the neutrino phase lattice) --- with
$\rho^4=-\rho$, so $\delta_{\mathrm{PMNS}}=\delta_{\mathrm{CKM}}^{\mathrm{lead}}+\pi$ is the quark phase times
the $\Z_2$ sheet ($\rho^3=-1$); the frame orientations of (iii) are correspondingly sheet-flipped
($\pm21\sin(\pi/3)$). So Target~D's CP sector is \emph{not} two independent phase inputs but \emph{one}
hexagonal unit read on the two sheets --- the magnitude residual $3\lambda_C^2$ (quark, \vref{v88}) and the
physical derivation stay \mC, but one of the two phase inputs is removed. \textbf{(v) The universal triality
phase (\veri{v233\_cp\_triality\_phase.py}).} In the canonical $\mu_6=\mu_3$(family/triality)$\times\mu_2$(sheet)
factorisation, $\rho=\omega^2(-1)$ with $\omega=e^{2\pi i/3}$ the $\Z_3$ triality centre (the $2/3$ cusp
weight); since $\rho^4=\rho^1\rho^3$ with $\rho^3\in\mu_2$, the quark ($\rho^1$) and lepton ($\rho^4$) phases
have the \emph{same} family/triality class and differ \emph{only} by the sheet. So the CP phase \emph{is} the
universal triality phase, sheet-split --- not even a free choice of $\mu_6$ power: the power choice of (iv)
is removed, and $\omega$ is the $\mu_3$ factor of the order-$30$ $(2,3,5)$ Coxeter/Milnor monodromy that
builds the hull (\vref{v232}). The physical CP derivation still stays \mC.
\end{winbox}

\begin{winbox}[Global look-elsewhere closure \mE\ (\veri{v100\_numerology\_null\_mc.py})]
The per-target look-elsewhere traps recorded above (the $\pi/3+2\lambda_C^2$ alternative in Target D, the
$\ln(46080)$ scale reading destroyed in \veri{v99\_koide\_flow\_time.py}) are local instances of the global
adversarial question: \emph{could the whole scorecard be formula-fishing?} The numerology null test answers
it with an explicit, declared null model run in adversarial mode: the \emph{entire} complexity-matched
formula grammar (which provably contains every scored TFPT formula --- the null space includes the theory)
is enumerated against conservative data windows. Joint formula-fishing probability $\prod_i p_i=10^{-25.8}$
over the $13$ scored observables; all $94\,500$ $F_{U(1)}$ variants solved, exactly one hits CODATA
($p_\alpha\le1.1{\cdot}10^{-5}$); total $\le10^{-30.7}$. The test is falsifiable in both directions: the
negative controls (seed perturbation $1\%/10\%$ collapses TFPT's own hit count $13\to9/6$; data shuffling
$\to1.2$) prove it has power, and the census is stable under budget/window variation ($<2$ orders). Honest
limit, stated as such everywhere: this is a null-model rejection \emph{conditional on the declared grammar}
--- no statistical test certifies a theory; the tension observables keep their large windows and contribute
nothing to the number. Details and the grammar definition: the audit note (\texttt{tfpt\_3}) and the module
docstring.
\end{winbox}

\begin{winbox}[External-review residual map \mE/\mC\ (\veri{v182\_reviewer\_residual\_map.py})]
A careful external review (2026-06-14) raised four concerns at the \mC/\mO\ seams: \textbf{(A)} the
mapping $2D\!\to\!4D$ (``why does the dynamics read the masses as the Pl\"ucker norm $11$?''),
\textbf{(B)} the abstractness of UV gravity, \textbf{(C)} the meta-look-elsewhere of the grammar choice,
and \textbf{(D)} the Koide M\"obius flow (non-integer flow time $t\!\sim\!2.84$, missing QFT generator).
None adds an \emph{independent} open gap. \textbf{A} \mE/\mC: the readout integer is exact ($11=
\sum_{k\le2}\binom4k=\dim S^+{-}\gcar$, the QBL boundary count, \vref{v174}) and the $2D\!\to\!4D$ step
is the holographic reduction, so $A=\texttt{QGEO.SYM.01}+F_{\mathrm{transfer}}$. \textbf{B} \mE:
``gravity is the geometric readout of the boundary'' \emph{is} \texttt{QGEO.SYM.01}. \textbf{D} \mE/\mC:
$\mathrm{Aut}(\PP)=\mathrm{PSL}(2,\mathbb C)$ is the M\"obius group and the cusp flow lives on the
deck $\PP\setminus\mu_4$, so the M\"obius \emph{nature} is forced by the geometry; the non-integer
$t\!\sim\!2.84$ only locates the physical point and the missing generator \emph{is} $F_{\mathrm{transfer}}$,
so $D=\texttt{QGEO.SYM.01}+F_{\mathrm{transfer}}$. So A, B, D collapse onto the \emph{two already-named}
residuals; only \textbf{C} is genuinely \emph{experiment-only}: the grammar is frozen (\vref{v84}),
minimal and over-determined (the anchor $a=(1,1,2)$ lands in $\ge6$ mutually-independent,
not-selected-for structures), a look-elsewhere-resistant signature whose residual --- correctly ---
falls only to out-of-sample data (JUNO, CMB-S4), never to argument. The four concerns are a unification
of the residual \emph{structure}, not a closure of the underlying premises, which stay \mO/\mC.
\end{winbox}

\end{document}
